% Section 18: Conclusion — reconstructed per CEA refinement protocol (2026-08-29)

\section{Conclusion}
\label{sec:conclusion}

We presented the Bounded-Authority Advisory Architecture (BAA), a decision-support framework that separates factual authority from generative language in agricultural AI. An 11-slot, single-record, fail-closed certification contract structurally prevents large language models from altering, hallucinating, or misbinding safety-critical chemical dosages, pre-harvest intervals, and application guidelines.

The architecture's central safety property was evaluated exhaustively: all 11 mutation classes across 40 base agronomic records (440 coverage cells) were evaluated against the certification predicate, which rejected every case. The B5-vs-B6 isolation demonstrates through controlled component ablation that slot completeness accounts for the corresponding corruption vulnerabilities. An empirical slot ablation confirms that every contract slot is individually load-bearing, with crop validation ($+$8.94\,pp) and provenance hash integrity ($+$8.92\,pp) creating the largest single-slot hazard surges across the benchmark when disabled. A benign-mutation control confirms the predicate discriminates: 0.0\% false rejection on approved chemicals, with a 3.67\% overall rejection rate attributable exclusively to conservative refusal of restricted compounds.

Measurement of existing baseline systems demonstrates the problem the architecture addresses: unconstrained LLMs accept unsafe recommendations in 8.0\% of live-API cases; LLM-judge guardrails perform worse under temporal conflict (75.0\% obsolete-gazette citation), not better. On a 100-case live benchmark, BAA achieves 83.0\% certified advisory correctness and reduces critical unsafe acceptance to 1.0\% (versus 15.0\% for the post-hoc LLM judge). The fail-closed design trades coverage for safety: 3.67\% of valid queries with banned or restricted chemical histories are conservatively abstained, a cost we report as a design parameter rather than a limitation to minimise.

Future work will extend this framework to field trials evaluating real crop yield preservation and applicator safety in partnership with national extension services.